\documentclass[conference]{IEEEtran}

\usepackage{cite}
\usepackage{amsmath,amssymb,amsfonts}
\usepackage{graphicx}
\usepackage{textcomp}
\usepackage{xcolor}
\usepackage{booktabs}
\usepackage{array}
\usepackage{url}
\usepackage[hidelinks]{hyperref}
\usepackage{microtype}
\usepackage{balance}

\newcommand{\figplaceholder}[2]{%
\begin{center}
\fbox{\begin{minipage}[c][#2][c]{0.92\linewidth}\centering #1\end{minipage}}
\end{center}}

\def\BibTeX{{\rm B\kern-.05em{\sc i\kern-.025em b}\kern-.08em
    T\kern-.1667em\lower.7ex\hbox{E}\kern-.125emX}}

\begin{document}

\title{ReSceneScribe: Static--Dynamic 3D Accident Reconstruction from Vehicle-Mounted Cameras and Multi-Agent Simulation for Evidence-Grounded Investigation}

\author{
\IEEEauthorblockN{Woosup Lee\IEEEauthorrefmark{1} and Yongmin Kim\IEEEauthorrefmark{2}}
\IEEEauthorblockA{\IEEEauthorrefmark{1}Planby Technologies Inc., Seoul, Republic of Korea}
\IEEEauthorblockA{\IEEEauthorrefmark{2}Department of Civil and Environmental Engineering, KAIST, Daejeon, Republic of Korea}
\IEEEauthorblockA{Corresponding author: Yongmin Kim, Email: yongmin@kaist.ac.kr}
}

\maketitle

\begin{abstract}
Traffic accident investigation requires metric information on vehicle positions,
road geometry, and collision relationships. However, accident scenes are often
cleared rapidly, forcing investigators to rely on partial videos captured by
vehicle-mounted cameras. This study presents ReSceneScribe, a static--dynamic
3D reconstruction framework that converts fragmented accident observations into
two practical outputs: an interpretable temporal replay and a collision-state
3D asset for evidence-grounded reporting. The proposed framework combines two
complementary data regimes. A pilot experiment using two real vehicles equipped
with windshield-mounted cameras verifies that ordinary vehicle-mounted sensing
can support static scene reconstruction under practical capture conditions.
Since repeated physical collision data collection is unsafe and inherently
limited, a four-vehicle scenario from DeepAccident Town03 is further used to
evaluate dynamic accident reconstruction in a controlled multi-agent
environment. The framework aligns multi-agent observations within a shared
metric coordinate system, removes dynamic objects from the static road
background, reconstructs vehicle geometry through vehicle-local motion
compensation, and constrains report generation using frame-indexed geometric
facts. In the selected scenario, the two involved vehicles overlap during source
frames 54--56, and the final collision-state frame shows an
oriented-bounding-box gap of $-0.4479$ m. The final hybrid Gaussian/RGB
reconstruction contains 1,177,009 vertices, comprising 1,100,000 static
background points and 77,009 motion-compensated vehicle points. These results
demonstrate that ReSceneScribe can support accident-scene reconstruction by
linking field-feasible sensing, simulation-based dynamic validation, and
auditable geometric reporting within a unified analysis framework.
\end{abstract}

\begin{IEEEkeywords}
traffic accident reconstruction, 3D reconstruction, vehicle-mounted camera,
static--dynamic decomposition, multi-agent sensing, Gaussian-compatible point
cloud, evidence-grounded reporting
\end{IEEEkeywords}

\section{Introduction}

Accident analysis depends on the spatial relationship among vehicles, lanes,
road markings, impact locations, and surrounding infrastructure. The most
reliable evidence is therefore not only a video record but a metric
three-dimensional state of the scene. In real accident response, however, this
state is rarely preserved for long. Vehicles and debris are cleared to reduce
secondary risk and to restore road capacity, while road-surface marks and
vehicle positions may be altered before detailed measurement is possible. As a
result, many cases are interpreted from partial photographs, vehicle-mounted
camera videos, and witness statements. These sources are useful, but a single
camera view is fundamentally limited by perspective distortion, occlusion, and
uncertain metric scale.

Conventional 3D accident documentation has used terrestrial laser scanning,
UAV photogrammetry, and related surveying tools to preserve scene geometry
\cite{cappelletti2019forensic,almeshal2020accuracy}. These methods can produce
accurate spatial records, but they require equipment, scene access, and
operator time that are not always available in routine accident handling.
Recent computer-vision methods provide a complementary direction.
Structure-from-Motion estimates camera motion and sparse geometry from image
correspondences \cite{schonberger2016sfm}; Neural Radiance Fields recover
continuous radiance and density fields from posed images \cite{mildenhall2020nerf};
and 3D Gaussian Splatting represents scenes with explicit Gaussian primitives
for efficient rendering \cite{kerbl2023gaussian}. Dynamic extensions such as 4D
Gaussian Splatting further model time-varying scenes \cite{wu2024fourDGS}.
However, direct application of these methods to accident analysis remains
insufficient because an accident scene is not merely a visual rendering target.
It must preserve which vehicle moved where, which road geometry was static,
when contact occurred, and which coordinates support each textual conclusion.

This paper focuses on an application-oriented reconstruction problem: how to
turn practical vehicle-mounted sensing and simulation-based accident data into
an interpretable and reusable accident reconstruction output. The key design
principle is static--dynamic separation. The road environment is reconstructed
as a stable static asset after removing moving-object traces, while vehicles
are represented separately by frame-wise poses and by motion-compensated
collision-state geometry. This prevents the common failure mode in which moving
vehicles are fused into the background and appear as repeated or blurred
shapes.

The study uses two complementary data sources. First, a physical pilot
experiment with two real vehicles and windshield-mounted cameras is used to
verify the practicality of vehicle-mounted acquisition for static road-scene
reconstruction. Second, a DeepAccident scenario generated in the CARLA simulator
is used for the dynamic collision stage \cite{dosovitskiy2017carla,wang2024deepaccident}.
DeepAccident is suitable here because it provides synchronized multi-view
camera, LiDAR, calibration, and label data from multiple vehicles, enabling
repeatable accident reconstruction without creating unsafe physical collisions.

The main contributions of this study are as follows: (1) a compact
static--dynamic decomposition pipeline for converting multi-agent accident
observations into a clean road background and dynamic vehicle trajectories; (2)
a vehicle-local motion compensation strategy for reconstructing
collision-state vehicle geometry while avoiding world-space ghosting artifacts;
(3) a dual-output accident-analysis application that provides an interactive
temporal replay for human interpretation and a Gaussian-compatible RGB PLY
asset for reusable 3D inspection; and (4) an evidence-grounded reporting
mechanism that constrains accident statements using vehicle identifiers, frame
indices, metric coordinates, and collision diagnostics.

\section{Method}

\subsection{Framework overview}

The proposed framework is organized around the distinction between what remains
fixed and what moves. Static elements include the road surface, lane markings,
buildings, roadside structures, and fixed infrastructure. Dynamic elements
include the target vehicles and other labelled moving objects. Treating all
observations as one static point set would maximize the number of points but
would destroy accident interpretability, because vehicle surfaces would be
accumulated at several positions. ReSceneScribe therefore constructs two linked
representations: a static background map and a dynamic vehicle representation.

\begin{figure*}[t!]
    \centering
    \IfFileExists{images/Figure1.png}{%
        \includegraphics[width=\textwidth]{images/Figure1.png}%
    }{%
        \figplaceholder{Workflow placeholder: field pilot reconstruction, simulation-based multi-vehicle accident replay, static--dynamic separation, vehicle-local fusion, and evidence-grounded reporting.}{28mm}%
    }
    \caption{Overview of ReSceneScribe. A physical pilot verifies practical vehicle-mounted acquisition, while the simulation scenario supplies repeatable dynamic accident observations.}
    \label{fig:workflow}
\end{figure*}

\subsection{Metric alignment of multi-agent observations}

For each vehicle agent $a$ and frame $t$, calibration provides an ego-to-world
transform $T^{w}_{a,t}$. A LiDAR point $p^{\ell}_{a,t}$ is mapped to the shared
world coordinate system by
\begin{equation}
    p^{w}_{a,t}=T^{w}_{a,t}T^{e}_{\ell,a,t}p^{\ell}_{a,t},
\end{equation}
where $T^{e}_{\ell,a,t}$ is the LiDAR-to-ego transform. Camera color is assigned
by projecting the world-space point into synchronized camera images through the
corresponding LiDAR-to-camera and intrinsic matrices. Points that are visible
in a valid camera image receive RGB values from the projected pixel, while
points without valid projection use an intensity-based fallback.

\subsection{Static background reconstruction}

The static background is generated by fusing points from all available agents
and frames after removing dynamic objects. Let $\mathcal{P}^{w}_{a,t}$ be the
set of world-space points at frame $t$ from agent $a$, and let
$\mathcal{B}_{d,t}$ be the set of 3D bounding boxes corresponding to dynamic
objects. A point is accepted into the static set only if it is outside all
dynamic boxes:
\begin{equation}
    \mathcal{S}=\{p \in \mathcal{P}^{w}_{a,t}: p \notin B,\; \forall B\in \mathcal{B}_{d,t}\}.
\end{equation}
This filtering also removes points inside the target vehicle boxes, including
self-vehicle observations, because target vehicles are reconstructed
separately. The filtered points are then voxel-downsampled to control density
while preserving road topology and large-scale structural features.

\subsection{Vehicle-local motion compensation}

Target vehicle geometry is reconstructed using the opposite strategy. Instead
of discarding points inside target vehicle boxes, these points are collected as
vehicle observations. For a target vehicle $k$, any point $p^{w}_{k,t}$ inside
its oriented bounding box at frame $t$ is transformed into the vehicle-local
coordinate system:
\begin{equation}
    p^{v}_{k,t}=\left(T^{w}_{k,t}\right)^{-1}p^{w}_{k,t}.
\end{equation}
Points are filtered to reduce road-surface contamination under the vehicle
body, accumulated across frames and observing agents, and voxel-downsampled in
local coordinates. The downsampled local geometry $\tilde{p}^{v}_{k}$ is then
transformed to the selected collision-state pose $T^{w}_{k,t_c}$:
\begin{equation}
    \hat{p}^{w}_{k}=T^{w}_{k,t_c}\tilde{p}^{v}_{k}.
\end{equation}

This procedure is motion compensation in the accident-reconstruction sense. If
vehicle points were accumulated directly in the world frame, the resulting
geometry would contain a trail along the vehicle trajectory. By accumulating in
the vehicle-local frame, each observed vehicle surface contributes to one
coherent vehicle geometry, which is then placed at the final collision pose.

\subsection{Collision diagnostics and grounded reporting}

The temporal replay uses frame-wise vehicle poses to animate all participating
vehicles. Collision validation is performed using oriented bounding boxes
(OBBs) projected to the ground plane. Let $g_{ij}(t)$ be the minimum
separating-axis gap between vehicles $i$ and $j$ at frame $t$. A positive value
indicates separation, while a negative value indicates overlap. The selected
collision-state frame is the last available frame in which the collision pair
is in a validated overlap state.

\section{Experiments and Results}

\subsection{Physical pilot experiment}

The physical pilot experiment was conducted to examine the feasibility of
vehicle-mounted image acquisition in a real road environment before relying on
a simulated collision dataset. Two real vehicles, a Hyundai Avante and a
Hyundai Casper, were used. Smartphone cameras were mounted near the windshield
area to approximate ordinary vehicle-mounted camera placement. The cameras
recorded RGB videos at a resolution of 1920$\times$1080 pixels and 30 frames/s.
The vehicles performed controlled approach and avoidance maneuvers in a
non-collision setting.

\begin{figure}[t]
    \centering
    \IfFileExists{images/Figure2.png}{%
        \includegraphics[width=\linewidth]{images/Figure2.png}%
    }{%
        \figplaceholder{Physical pilot placeholder: two real vehicles and static reconstruction from vehicle-mounted RGB video.}{30mm}%
    }
    \caption{Physical pilot experiment and static scene reconstruction using vehicle-mounted RGB cameras.}
    \label{fig:physical_pilot_experiment}
\end{figure}

\subsection{DeepAccident simulation scenario}

The simulation experiment uses the DeepAccident mini/local scenario. The
scenario is a four-way junction under \texttt{MidRainSunset} conditions. It
contains four vehicle agents: two collision vehicles and two following
vehicles. Each vehicle provides 56 available frames with multi-view RGB
cameras, LiDAR, calibration files, and 3D object labels. The source frame range
used in this study is 001--056.

\begin{table}[t]
\centering
\footnotesize
\caption{Essential experimental configuration.}
\label{tab:setup}
\begin{tabular}{p{0.36\linewidth}p{0.56\linewidth}}
\toprule
Item & Configuration \\
\midrule
Physical pilot & Two real vehicles; windshield-mounted smartphone cameras; 1920$\times$1080; 30 frames/s; controlled non-collision approach and avoidance maneuvers \\
Simulation scenario & DeepAccident Town03, type001/subtype0001/scenario00024 \\
Scene condition & Four-way junction, \texttt{MidRainSunset} \\
Simulation agents & Four vehicle agents: two collision vehicles and two following vehicles \\
Available source frames & 001--056 \\
Main outputs & Interactive temporal replay; collision-state hybrid Gaussian/RGB PLY \\
\bottomrule
\end{tabular}
\end{table}

\subsection{Temporal accident replay}

The first output is an interactive four-vehicle accident replay. The static
background is reconstructed by fusing LiDAR streams from all four vehicle
agents after dynamic-object removal. Vehicle motion is replayed using
calibration-derived frame-wise poses and dimension-aware vehicle proxies. This
output is intended for human interpretation: a user can inspect the road
layout, vehicle trajectories, and synchronized vehicle-mounted camera frames in
one view.

\begin{figure*}[t!]
    \centering
    \IfFileExists{images/Figure3.png}{%
        \includegraphics[width=\textwidth]{images/Figure3.png}%
    }{%
        \figplaceholder{Result placeholder: interactive replay with fused static background, four vehicle trajectories, synchronized vehicle-mounted camera panels, and collision diagnostics.}{28mm}%
    }
    \caption{Simulation-based accident replay from the DeepAccident Town03 four-vehicle collision scenario.}
\label{fig:replay}
\end{figure*}

The four vehicle trajectories verify that the scenario is not a single-view
reconstruction problem. The two collision vehicles approach the intersection
from different directions, while the following vehicles provide additional
viewpoints and context. In the selected frame range, the collision vehicles
overlap from frame 54 to frame 56. At frame 56, the final OBB gap is
$-0.4479$ m and the center distance is 4.2299 m, supporting the use of frame 56
as the final collision-state frame.

\subsection{Collision-state hybrid Gaussian/RGB reconstruction}

The second output is a collision-state 3D asset. The method starts from
8,616,764 source LiDAR points collected across the four vehicle agents. During
static background generation, 312,543 dynamic or target-vehicle points are
removed to suppress moving-object contamination. The final PLY contains
1,177,009 vertices: 1,100,000 static background points and 77,009 vehicle
points reconstructed through vehicle-local motion compensation. The PLY
includes ordinary RGB fields for general point-cloud viewers and
Gaussian-compatible fields such as DC color, opacity, scale, and rotation for
downstream Gaussian-style rendering workflows.

\begin{table}[t]
\centering
\footnotesize
\caption{Vehicle-level motion and reconstructed geometry summary.}
\label{tab:vehicle}
\begin{tabular}{lrr}
\toprule
Vehicle role & Net motion (m) & Final points \\
\midrule
Collision vehicle A & 39.9009 & 28,979 \\
Follower behind A & 32.6585 & 8,024 \\
Collision vehicle B & 43.8988 & 28,889 \\
Follower behind B & 41.2916 & 11,117 \\
\bottomrule
\end{tabular}
\end{table}

\begin{table}[t]
\centering
\footnotesize
\caption{Core quantitative results used for accident interpretation.}
\label{tab:results}
\begin{tabular}{lr}
\toprule
Metric & Value \\
\midrule
Verified OBB overlap frames & 54--56 \\
Selected collision-state frame & 56 \\
Final OBB gap & $-0.4479$ m \\
Final center distance & 4.2299 m \\
Source LiDAR points & 8,616,764 \\
Removed dynamic/target-vehicle points & 312,543 \\
Final PLY vertices & 1,177,009 \\
Static / vehicle PLY points & 1,100,000 / 77,009 \\
\bottomrule
\end{tabular}
\end{table}

\section{Discussion}

The proposed framework is intentionally different from a general-purpose visual
reconstruction pipeline. General neural rendering methods often prioritize
photorealistic novel-view synthesis. Accident reconstruction additionally
requires metric traceability. A visually plausible reconstruction is not enough
if the report cannot identify which frame, coordinate, or vehicle measurement
supports a claim. By maintaining a static background, dynamic vehicle
trajectories, collision diagnostics, and structured reporting facts,
ReSceneScribe keeps the visual, geometric, and linguistic components tied to
the same evidence base.

Several limitations remain. First, the final PLY is a deterministic hybrid
Gaussian/RGB representation rather than a fully optimized dynamic 3D Gaussian
Splatting model. It is designed for metric stability and viewer compatibility,
not for maximum photorealism. Second, vehicle surface completeness depends on
LiDAR visibility; following vehicles or heavily occluded surfaces may have
lower point density. Third, collision validation uses OBB overlap as a
geometric proxy. This is appropriate for scenario-level accident-state
verification, but future deployment should incorporate deformation, contact
surface estimation, and uncertainty propagation.

\section{Conclusion}

This paper presented ReSceneScribe, a static--dynamic 3D reconstruction
framework for multi-vehicle accident replay and evidence-grounded reporting.
The work contributes an application-oriented pipeline that connects a
real-world vehicle-mounted camera pilot with simulation-based dynamic collision
reconstruction. The method separates static road geometry from dynamic
vehicles, reconstructs vehicles using vehicle-local motion compensation, and
exports both an interactive replay and a collision-state hybrid Gaussian/RGB
PLY. In the selected scenario, the final asset contains 1,177,009 vertices and
the collision pair is validated by a final OBB gap of $-0.4479$ m at frame 56.

\section*{Code and Data Availability}

The implementation of ReSceneScribe is available at:
\url{https://github.com/WoosopYi/ReSceneScribe.git}. Raw DeepAccident data are
not redistributed in the repository and must be obtained from the dataset
provider.

\section*{Declaration of Generative AI and AI-Assisted Technologies}

Generative AI assistance was used only to support visual drafting and layout
refinement of research figures. The authors reviewed and finalized all
scientific content and take full responsibility for the article.

\balance

\begin{thebibliography}{99}

\bibitem{cappelletti2019forensic}
C.~Cappelletti, M.~Boniardi, A.~Casaroli, C.~I. De~Gaetani, D.~Passoni,
L.~Pinto, ``Forensic engineering surveys with UAV photogrammetry and laser
scanning techniques,'' \emph{Int. Arch. Photogramm. Remote Sens. Spatial Inf.
Sci.}, vol. XLII-2/W9, pp. 227--234, 2019.

\bibitem{almeshal2020accuracy}
A.~M. Almeshal, M.~R. Alenezi, and A.~K. Alshatti, ``Accuracy assessment of
small unmanned aerial vehicle for traffic accident photogrammetry in the
extreme operating conditions of Kuwait,'' \emph{Information}, vol. 11, no. 9,
p. 442, 2020.

\bibitem{schonberger2016sfm}
J.~L. Sch\"{o}nberger and J.-M. Frahm, ``Structure-from-motion revisited,'' in
\emph{Proc. IEEE Conf. Computer Vision and Pattern Recognition}, 2016,
pp. 4104--4113.

\bibitem{mildenhall2020nerf}
B.~Mildenhall, P.~P. Srinivasan, M.~Tancik, J.~T. Barron, R.~Ramamoorthi, and
R.~Ng, ``NeRF: Representing scenes as neural radiance fields for view
synthesis,'' in \emph{Proc. European Conf. Computer Vision}, 2020,
pp. 405--421.

\bibitem{kerbl2023gaussian}
B.~Kerbl, G.~Kopanas, T.~Leimk\"{u}hler, and G.~Drettakis, ``3D Gaussian
Splatting for real-time radiance field rendering,'' \emph{ACM Trans. Graph.},
vol. 42, no. 4, pp. 1--14, 2023.

\bibitem{wu2024fourDGS}
G.~Wu, T.~Yi, J.~Fang, L.~Xie, X.~Zhang, W.~Wei, W.~Liu, Q.~Tian, and X.~Wang,
``4D Gaussian Splatting for real-time dynamic scene rendering,'' in
\emph{Proc. IEEE/CVF Conf. Computer Vision and Pattern Recognition}, 2024,
pp. 20310--20320.

\bibitem{dosovitskiy2017carla}
A.~Dosovitskiy, G.~Ros, F.~Codevilla, A.~Lopez, and V.~Koltun, ``CARLA: An
open urban driving simulator,'' in \emph{Proc. 1st Annual Conf. Robot
Learning}, PMLR 78, 2017, pp. 1--16.

\bibitem{wang2024deepaccident}
T.~Wang, S.~Kim, J.~Wenxuan, E.~Xie, C.~Ge, J.~Chen, Z.~Li, and P.~Luo,
``DeepAccident: A motion and accident prediction benchmark for V2X autonomous
driving,'' \emph{Proc. AAAI Conf. Artif. Intell.}, vol. 38, no. 6,
pp. 5599--5606, 2024.

\end{thebibliography}

\end{document}
